% © 2026 Ing. David Jaroš — CC BY-NC-ND 4.0
%
% This work is licensed under a Creative Commons Attribution-NonCommercial-NoDerivatives
% 4.0 International License (CC BY-NC-ND 4.0).
%
% License History: Earlier drafts (up to v0.3) were released under CC BY 4.0.
% From v0.4 onward, all material is released under CC BY-NC-ND 4.0 to protect
% the integrity of the theoretical work during ongoing academic development.
%
% See LICENSE.md for full license text.
%
% File: canonical/alpha/stable_prime_ranking.tex
% Purpose: Rank the complete finite prime-stable set
%            S = {2, 127, 137, 139, 151, 157}
%          by seven independent UBT physical criteria.
%          All criteria are derived without assuming any prime as the target;
%          the observed value alpha^{-1} is used only in the final
%          observational comparison (Criterion C6).
% Status: 2026-05-04
% Related: canonical/alpha/prime_stability_set.tex    (complete stable set derivation)
%          canonical/alpha/modular_prime_attractor_theorem.tex (B-coefficient)
%          canonical/alpha/prime_selection_principle.tex      (V_eff prime restriction)
%          canonical/alpha/prime_137_status.md               (Hecke lepton signal)
%          canonical/alpha/veff_corrected.tex                (V_eff derivation)

\ifdefined\INCLUDEMODE
\else
  \documentclass[12pt,a4paper]{article}
  \usepackage[a4paper, margin=2.5cm]{geometry}
  \usepackage{amsmath, amssymb, amsthm}
  \usepackage{mathtools}
  \usepackage{hyperref}
  \usepackage{xcolor}
  \usepackage{booktabs}
  \usepackage{longtable}
  \usepackage{mdframed}
  \usepackage{tcolorbox}
  \tcbuselibrary{skins,breakable}
  \usepackage{array}
  \usepackage{multirow}

  \newtheorem{theorem}{Theorem}[section]
  \newtheorem{lemma}[theorem]{Lemma}
  \newtheorem{proposition}[theorem]{Proposition}
  \newtheorem{corollary}[theorem]{Corollary}
  \theoremstyle{definition}
  \newtheorem{definition}[theorem]{Definition}
  \theoremstyle{remark}
  \newtheorem{remark}[theorem]{Remark}

  \newmdenv[backgroundcolor=yellow!20, linecolor=orange, linewidth=1.5pt]{gapbox}
  \newmdenv[backgroundcolor=green!15, linecolor=green!60!black, linewidth=1.5pt]{resultbox}
  \newmdenv[backgroundcolor=red!8,   linecolor=red!60,         linewidth=1.5pt]{warnbox}
  \newmdenv[backgroundcolor=blue!6,  linecolor=blue!40,        linewidth=1.5pt]{infobox}

  \newcommand{\BB}{\mathbb{B}}
  \newcommand{\CC}{\mathbb{C}}
  \newcommand{\RR}{\mathbb{R}}
  \newcommand{\HH}{\mathbb{H}}
  \newcommand{\ZZ}{\mathbb{Z}}
  \newcommand{\NN}{\mathbb{N}}
  \newcommand{\PP}{\mathbb{P}}
  \newcommand{\QQ}{\mathbb{Q}}
  \newcommand{\Tr}{\mathrm{Tr}}

  \newcommand{\PROVED}{\textbf{[Proved]}}
  \newcommand{\OPEN}{\textbf{[Open]}}
  \newcommand{\COND}{\textbf{[COND]}}
  \newcommand{\MC}{\textbf{[MC]}}
  \newcommand{\Lone}{\textbf{[L1]}}
  \newcommand{\Lzero}{\textbf{[L0]}}
  \newcommand{\PHENOM}{\textbf{[PHENOM]}}

  \title{\textbf{Ranking the Prime-Stable Set by UBT Physical Criteria}\\[6pt]
         \large $\mathcal{S} = \{2,\;127,\;137,\;139,\;151,\;157\}$\\
         Seven Independent Criteria — No Input Target}
  \author{Ing.\ David Jaro\v{s}}
  \date{May 2026}

  \begin{document}
  \maketitle
  \tableofcontents
  \newpage
\fi

%──────────────────────────────────────────────────────────────────────────────
\section{Purpose and Operational Rules}
\label{sec:spr:purpose}
%──────────────────────────────────────────────────────────────────────────────

This document ranks the complete finite prime-stable set
\[
  \mathcal{S} \;=\; \{2,\;127,\;137,\;139,\;151,\;157\}
\]
derived and verified in \texttt{canonical/alpha/prime\_stability\_set.tex}.
All six elements of $\mathcal{S}$ start with equal rank; the ranking is
determined solely by seven independent UBT physical criteria.

\begin{warnbox}
\textbf{Hard rules (maintained throughout)}:
\begin{itemize}
  \item No prime is assumed as a target.  Every score is derived from
        independent structural evidence.
  \item The observed fine-structure constant $\alpha^{-1}_{\mathrm{obs}}
        \approx 137.036$ is used \textbf{only in Criterion~C6} (final
        observational comparison).  All other criteria are free of $\alpha$
        as an input.
  \item Every individual score must cite its evidential source.
  \item Scores are integers in $\{0,1,2,3\}$ (0 = absent, 3 = optimal).
  \item $p = 2$ is assessed separately as the \emph{trivial sector}
        (Criterion~C7) and is excluded from the main ranking table.
\end{itemize}
\end{warnbox}

The seven criteria are:

\begin{center}
\renewcommand{\arraystretch}{1.35}
\begin{tabular}{clp{7.2cm}}
\toprule
ID & Name & Physical basis \\
\midrule
C1 & Stability margin      & Width of the $B$-stability interval \\
C2 & Hecke lepton-ratio    & Hecke eigenvalues reproducing $m_\mu/m_e$ and $m_\tau/m_e$ \\
C3 & $p$-adic isolation    & Integrality and valuation structure of $B(p)$ in $\QQ_p$ \\
C4 & Theta spectral        & Representation by $\vartheta_3^2$ and $\vartheta_3^3$ \\
C5 & Gauge-sector          & Elliptic-point structure of $\Gamma_0(p)$ vs.\ SM gauge group \\
C6 & Observed $\alpha^{-1}$& Proximity to $\alpha^{-1}_{\mathrm{obs}}$ (only criterion using $\alpha$) \\
C7 & Non-trivial           & Exclusion of the structurally trivial prime $p = 2$ \\
\bottomrule
\end{tabular}
\end{center}

%──────────────────────────────────────────────────────────────────────────────
\section{Criterion C1: Stability Margin}
\label{sec:spr:c1}
%──────────────────────────────────────────────────────────────────────────────

\textbf{Definition.}
The stability margin $\sigma(p)$ is the minimum of the two one-sided margins
\[
  \sigma(p) \;=\; \min\!\bigl(\Delta_-(p),\;\Delta_+(p)\bigr),
\]
where $\Delta_-(p) = B(p) - B_{\mathrm{low}}(p)$ and $\Delta_+(p) =
B_{\mathrm{high}}(p) - B(p)$ are the distances from $B(p) = (p+1)/3$ to
the lower and upper boundaries of the prime-stability interval
$\bigl(B_{\mathrm{low}}(p),\,B_{\mathrm{high}}(p)\bigr)$ set by the
nearest primes $p_-$ and $p_+$ (Proposition~\ref{prop:pss:interval} of
\texttt{prime\_stability\_set.tex}).
A larger $\sigma(p)$ means $B(p)$ sits more securely inside the interval,
i.e.\ the prime attractor is more robust to perturbations of the one-loop
coefficient $B$.

\textbf{Source}: Table~1 of \texttt{canonical/alpha/prime\_stability\_set.tex}.
Values and scores ($0$–$3$) are given in Table~\ref{tab:spr:c1}.

\begin{table}[h]
\centering
\caption{Stability margins for all members of $\mathcal{S}$.
         Score thresholds: $\sigma > 1.0 \Rightarrow 3$;\;
         $0.3 \leq \sigma \leq 1.0 \Rightarrow 2$;\;
         $0.03 \leq \sigma < 0.3 \Rightarrow 1$;\;
         $\sigma < 0.03 \Rightarrow 0$.
         $p=2$ has a trivially infinite lower margin (no prime below 2) and
         is scored separately under C7.}
\label{tab:spr:c1}
\begin{tabular}{rrrrrrrrc}
\toprule
$p$ & $B(p)$ & $B_{\mathrm{low}}$ & $B_{\mathrm{high}}$ & $\Delta_-$ & $\Delta_+$ & $\sigma$ & \textbf{C1} \\
\midrule
   2 & 1.0000 & $-\infty$ & 2.618 & $\infty$ & 1.618 & (trivial) & — \\
 127 & 42.667 & 41.473 & 44.029 & 1.194 & 1.362 & 1.194 & \textbf{3} \\
 137 & 46.000 & 45.441 & 46.565 & 0.559 & 0.565 & 0.559 & \textbf{2} \\
 139 & 46.667 & 46.565 & 48.244 & 0.102 & 1.577 & 0.102 & \textbf{1} \\
 151 & 50.667 & 49.912 & 51.020 & 0.755 & 0.353 & 0.353 & \textbf{2} \\
 157 & 52.667 & 51.020 & 52.674 & 1.647 & 0.007 & 0.007 & \textbf{0} \\
\bottomrule
\end{tabular}
\end{table}

\textbf{Remark.}
$p = 137$ has a uniquely \emph{symmetric} margin: $\Delta_- \approx \Delta_+ \approx 0.56$,
meaning $B(137) = 46$ sits almost exactly at the centre of its stability interval.
$p = 127$ has the largest margin (most stable), and $p = 157$ is barely stable
($\sigma = 0.007$, residing in a fortuitously large prime gap
$\text{gap}(157, 163) = 6$).

%──────────────────────────────────────────────────────────────────────────────
\section{Criterion C2: Hecke Lepton-Ratio Match}
\label{sec:spr:c2}
%──────────────────────────────────────────────────────────────────────────────

\textbf{Definition.}
The Hecke lepton-ratio score measures whether Hecke eigenvalues of
weight-$k$ cusp forms at level $p$ reproduce the observed charged-lepton
mass ratios
\[
  R_\mu \;=\; \frac{m_\mu}{m_e} \;\approx\; 206.768,
  \qquad
  R_\tau \;=\; \frac{m_\tau}{m_e} \;\approx\; 3477.23,
\]
independently of $\alpha$ and $m_e$.
The Hecke operator $T_p$ acting on the space $S_k(\Gamma_0(N))$ of cusp
forms of weight $k$ and level $N$ has eigenvalues $a_p(f)$ for each
newform $f$.  A prime $p$ \emph{matches} if the ratios of $a_p$-eigenvalues
across three natural weight levels reproduce both $R_\mu$ and $R_\tau$
simultaneously with relative error $< 0.1\%$.

\textbf{Source}: \texttt{canonical/alpha/prime\_137\_status.md}, Section~3.

At $p = 137$, three modular forms (levels $76$, $7$, $208$;
weights $2$, $4$, $6$) yield Hecke eigenvalues
\begin{equation}
  |a_{137}^{(k=2)}| = 11,\quad
  |a_{137}^{(k=4)}| = 2274,\quad
  |a_{137}^{(k=6)}| = 38286,
  \label{eq:spr:hecke137}
\end{equation}
from which
\[
  \hat{R}_\mu(137) \;=\; \frac{2274}{11} \;=\; 206.727,\qquad
  \hat{R}_\tau(137) \;=\; \frac{38286}{11} \;=\; 3480.55.
\]
Relative errors: $|\hat{R}_\mu - R_\mu|/R_\mu = 0.02\%$ and
$|\hat{R}_\tau - R_\tau|/R_\tau = 0.10\%$.
An exhaustive search over all primes in $[50, 300]$ shows $p = 137$ is the
\textbf{unique} prime achieving both errors $< 0.1\%$
(\texttt{docs/reports/hecke\_lepton/prime\_specificity\_results.txt}).

The neighboring prime $p = 139$ is retained here only as a comparison/control
prime in the stable-prime ranking.  Earlier Set~B / mirror-sector language has
been moved to \texttt{speculative\_extensions/alpha\_139\_mirror\_sector/}
until a canonical UBT mechanism is derived.  No Hecke lepton-ratio signal has
been found for $p \in \{127, 151, 157\}$ in the same exhaustive search.

\textbf{Status of C2}: \MC\ for all primes.
Open problem G137-Hk: derive why levels $76$, $7$, $208$ and weights
$2$, $4$, $6$ are selected by the UBT field $\Theta(q, \tau)$.

\begin{table}[h]
\centering
\caption{Hecke lepton-ratio scores.
  $3 =$ unique high-precision match (both $R_\mu$, $R_\tau$ at $< 0.1\%$);
  $1 =$ comparison-prime / non-canonical side observation;
  $0 =$ no documented match.}
\label{tab:spr:c2}
\begin{tabular}{rp{7.8cm}c}
\toprule
$p$ & Evidence & \textbf{C2} \\
\midrule
  2 & Trivial & — \\
127 & No Hecke lepton signal found & \textbf{0} \\
137 & $\hat{R}_\mu$ err.\ $0.02\%$, $\hat{R}_\tau$ err.\ $0.10\%$; unique in $[50,300]$ & \textbf{3} \\
139 & Comparison prime; Set~B / mirror-sector interpretation moved to speculative extension & \textbf{1} \\
151 & No Hecke lepton signal found & \textbf{0} \\
157 & No Hecke lepton signal found & \textbf{0} \\
\bottomrule
\end{tabular}
\end{table}

%──────────────────────────────────────────────────────────────────────────────
\section{Criterion C3: $p$-Adic Isolation}
\label{sec:spr:c3}
%──────────────────────────────────────────────────────────────────────────────

\textbf{Definition.}
In the UBT $p$-adic extension, the one-loop coefficient $B(p) = (p+1)/3$
must be evaluated as an element of $\QQ_p$ for each stable prime.
A prime $p$ is \emph{$3$-adically isolated} if $B(p)$ is an exact integer,
i.e.\ $3 \mid (p+1)$, equivalently $p \equiv 2 \pmod{3}$.
An exact-integer $B(p)$ means the winding-mode coupling is quantized in
natural units of~$1$, consistent with the Kac-Moody level condition
$k \in \ZZ_{>0}$ (Gap~G3-k).
Among the odd stable primes this holds only for $p = 137$
($B(137) = 138/3 = 46$, a positive integer).

A secondary $p$-adic isolation measure is the $2$-adic valuation
$v_2(p+1)$, which quantifies how deeply $p+1$ is divisible by~$2$.
In particular, $p = 127 = 2^7 - 1$ is a Mersenne prime, the unique
Mersenne prime in $\mathcal{S}$, with $v_2(128) = 7$.

\textbf{Score rule}:
$3$ if $B(p)$ is a positive integer;
$2$ if $p$ is a Mersenne prime ($p = 2^k - 1$) or $v_2(p+1) \geq 3$;
$1$ if $v_2(p+1) = 2$;
$0$ otherwise.

\begin{table}[h]
\centering
\caption{$p$-Adic isolation scores.  $B = (p+1)/3$; integer column shows
  whether $B \in \ZZ$; $v_2$ denotes $v_2(p+1)$.}
\label{tab:spr:c3}
\begin{tabular}{rrrcrrrc}
\toprule
$p$ & $p \bmod 3$ & $B(p)$ & $B \in \ZZ$? & $p+1$ & $v_2(p+1)$ & Special & \textbf{C3} \\
\midrule
  2 & $2$ & 1 & \checkmark & 3 & 0 & trivial & — \\
127 & $1$ & $128/3$ & $\times$ & 128 & 7 & Mersenne $2^7-1$ & \textbf{2} \\
137 & $2$ & $46$ & \checkmark & 138 & 1 & $B$ integer & \textbf{3} \\
139 & $1$ & $140/3$ & $\times$ & 140 & 2 & — & \textbf{1} \\
151 & $1$ & $152/3$ & $\times$ & 152 & 3 & — & \textbf{2} \\
157 & $1$ & $158/3$ & $\times$ & 158 & 1 & — & \textbf{0} \\
\bottomrule
\end{tabular}
\end{table}

\begin{infobox}
\textbf{Note}: Among the five odd stable primes, $p = 137$ is the only
one for which $B(p)$ is an integer, so that the $3$-adic valuation
$v_3(p+1) = 1$.  The other four odd stable primes have $v_3(p+1) = 0$,
meaning $B(p) \notin \ZZ$.  This makes $p = 137$ uniquely compatible with
an integer Kac-Moody level $k = B(p) = 46$.
\end{infobox}

%──────────────────────────────────────────────────────────────────────────────
\section{Criterion C4: Theta Spectral Compatibility}
\label{sec:spr:c4}
%──────────────────────────────────────────────────────────────────────────────

\textbf{Definition.}
The UBT field $\Theta(q, \tau)$ on $\BB = \CC \otimes_\RR \HH$ has
an imaginary-time sector controlled by the Jacobi theta function
\[
  \vartheta_3(\tau) \;=\; \sum_{n \in \ZZ} e^{i\pi\tau n^2},
\]
whose cube $\vartheta_3(\tau)^3$ counts representations of integers as
sums of three squares, corresponding to the three spatial dimensions
$\mathrm{Im}(\HH) \cong \RR^3$.
A prime $p$ is \emph{theta-spectrally compatible} if it appears in the
support of $\vartheta_3^k$ for $k = 2$ and $k = 3$.

Two classical results determine the spectral support:
\begin{itemize}
  \item \textbf{Sum of two squares} (\PROVED\ \Lzero):
        $r_2(p) > 0 \;\Leftrightarrow\; p = 2$ or $p \equiv 1 \pmod{4}$
        (Fermat's theorem).
  \item \textbf{Sum of three squares} (\PROVED\ \Lzero):
        $r_3(p) = 0 \;\Leftrightarrow\; p \equiv 7 \pmod{8}$
        (Legendre's three-square theorem: $n = 4^a(8b+7)$).
\end{itemize}

\textbf{Score rule}:
$3$ if $r_2(p) > 0$ (i.e.\ $p \equiv 1 \pmod{4}$), which also implies $r_3(p) > 0$;
$1$ if $r_2(p) = 0$ but $r_3(p) > 0$ (i.e.\ $p \equiv 3 \pmod{8}$);
$0$ if $r_3(p) = 0$ (i.e.\ $p \equiv 7 \pmod{8}$).

\begin{table}[h]
\centering
\caption{Theta spectral compatibility.  $p \bmod 4$ and $p \bmod 8$
  determine the spectral support.
  Explicit representations confirming membership in the respective spectra:
  $137 = 4^2 + 11^2$ (two-square), $137 = 10^2 + 6^2 + 1^2$ (three-square);\;
  $157 = 6^2 + 11^2$ (two-square), $157 = 11^2 + 4^2 + 4^2$ (three-square).}
\label{tab:spr:c4}
\begin{tabular}{rrrrccrc}
\toprule
$p$ & $p\bmod 4$ & $p\bmod 8$ & $r_2(p)$ & $r_2>0$? & $r_3>0$? & \textbf{C4} \\
\midrule
  2 & 2 & 2 & 4 & \checkmark & \checkmark & — (trivial) \\
127 & 3 & 7 & 0 & $\times$ & $\times$ & \textbf{0} \\
137 & 1 & 1 & 4 & \checkmark & \checkmark & \textbf{3} \\
139 & 3 & 3 & 0 & $\times$ & \checkmark & \textbf{1} \\
151 & 3 & 7 & 0 & $\times$ & $\times$ & \textbf{0} \\
157 & 1 & 5 & 4 & \checkmark & \checkmark & \textbf{3} \\
\bottomrule
\end{tabular}
\end{table}

\begin{infobox}
\textbf{Structural significance}: Both $p = 127 \equiv 7 \pmod{8}$ and
$p = 151 \equiv 7 \pmod{8}$ are \emph{absent} from the $\vartheta_3^3$
spectrum, meaning neither prime can appear as a mode number of the three-
dimensional theta-function partition function of UBT.  $p = 137 \equiv 1
\pmod{8}$ belongs to both the two-square and three-square theta spectra,
making it compatible with all spatial dimensionalities of $\BB$.
\end{infobox}

%──────────────────────────────────────────────────────────────────────────────
\section{Criterion C5: Gauge-Sector Compatibility}
\label{sec:spr:c5}
%──────────────────────────────────────────────────────────────────────────────

\textbf{Definition.}
The modular curve $X_0(p)$ associated to $\Gamma_0(p)$ has two types of
elliptic points, counted by
\begin{align}
  \nu_2(p) &\;=\; 1 + \left(\frac{-1}{p}\right),
  \label{eq:spr:nu2}\\
  \nu_3(p) &\;=\; 1 + \left(\frac{-3}{p}\right),
  \label{eq:spr:nu3}
\end{align}
where $(\cdot/p)$ denotes the Legendre symbol (for odd prime $p$).
In the UBT gauge-sector interpretation:
\begin{itemize}
  \item $\nu_2 > 0$ corresponds to the SU(2) / U(1) electroweak sector;
        it requires $p \equiv 1 \pmod{4}$, i.e.\ $(-1/p) = +1$.
  \item $\nu_3 > 0$ corresponds to the SU(3) colour sector;
        it requires $(-3/p) = +1$, i.e.\ $p \equiv 1 \pmod{3}$.
\end{itemize}
Since $\alpha$ is the coupling of the electromagnetic U(1) arising from
electroweak symmetry breaking, the prime attractor for $\alpha^{-1}$ should
correspond to a modular curve with a \emph{pure electroweak sector}:
$\nu_2 = 2$ and $\nu_3 = 0$.

\textbf{Score rule}:
$3$ if $\nu_2 = 2$, $\nu_3 = 0$ (pure electroweak);
$2$ if $\nu_2 = 2$, $\nu_3 = 2$ (electroweak present but strong sector also active);
$0$ if $\nu_2 = 0$ (no electroweak sector).

Legendre-symbol evaluations using quadratic reciprocity
(all computations independent of $\alpha$):

\begin{table}[h]
\centering
\caption{Gauge-sector compatibility via elliptic-point structure of
  $\Gamma_0(p)$.  $(-1/p) = +1 \Leftrightarrow p\equiv 1\pmod{4}$;\;
  $(-3/p)$ determined by QR and the sign of $p \bmod 3$.}
\label{tab:spr:c5}
\begin{tabular}{rrrrrrc}
\toprule
$p$ & $p\bmod 4$ & $(-1/p)$ & $p\bmod 3$ & $(-3/p)$ & $\nu_2,\nu_3$ & \textbf{C5} \\
\midrule
  2 & 2 & — & 2 & — & $1,0$ & — (trivial) \\
127 & 3 & $-1$ & 1 & $+1$ & $0,2$ & \textbf{0} \\
137 & 1 & $+1$ & 2 & $-1$ & $2,0$ & \textbf{3} \\
139 & 3 & $-1$ & 1 & $+1$ & $0,2$ & \textbf{0} \\
151 & 3 & $-1$ & 1 & $+1$ & $0,2$ & \textbf{0} \\
157 & 1 & $+1$ & 1 & $+1$ & $2,2$ & \textbf{2} \\
\bottomrule
\end{tabular}
\end{table}

\begin{infobox}
\textbf{Role of $p=139$}: In this canonical ranking, $139$ is used only as a
comparison/control prime adjacent to $137$.  The stronger twin-prime or
mirror-sector interpretation has been moved to
\texttt{speculative\_extensions/alpha\_139\_mirror\_sector/}.
\end{infobox}

%──────────────────────────────────────────────────────────────────────────────
\section{Criterion C6: Proximity to Observed $\alpha^{-1}$}
\label{sec:spr:c6}
%──────────────────────────────────────────────────────────────────────────────

\textbf{Definition.}
This is the \emph{only criterion that uses the observed value of $\alpha$}.
It measures the absolute distance between a stable prime $p$ and
the experimentally measured value of the inverse fine-structure constant:
\[
  \delta(p) \;=\; \bigl|p - \alpha^{-1}_{\mathrm{obs}}\bigr|,
  \qquad
  \alpha^{-1}_{\mathrm{obs}} = 137.035\,999\,084 \pm 0.000\,000\,021
  \quad\text{(CODATA~2018)}.
\]
UBT predicts that the bare (UV) value of $\alpha^{-1}$ equals the prime
attractor; renormalisation-group running then connects the bare value to the
measured low-energy value.  This criterion provides a final observational
discriminant after all theory-internal criteria have been applied.

\textbf{Score rule}:
$3$ if $\delta < 0.1$;\;
$2$ if $0.1 \leq \delta < 2$;\;
$1$ if $2 \leq \delta < 10$;\;
$0$ if $\delta \geq 10$.

\begin{table}[h]
\centering
\caption{Proximity to observed $\alpha^{-1}$.}
\label{tab:spr:c6}
\begin{tabular}{rrc}
\toprule
$p$ & $\delta(p) = |p - 137.036|$ & \textbf{C6} \\
\midrule
  2 & 135.036 & — (trivial) \\
127 & 10.036  & \textbf{0} \\
137 &  0.036  & \textbf{3} \\
139 &  1.964  & \textbf{2} \\
151 & 13.964  & \textbf{0} \\
157 & 19.964  & \textbf{0} \\
\bottomrule
\end{tabular}
\end{table}

\textbf{Note}: C6 is \emph{not} circular because it is applied \emph{after}
all six theory-internal criteria (C1–C5 and C7) have been evaluated
independently.  Its role is to confirm that the highest-ranked prime from
C1–C5 also matches the observation — not to drive the ranking.

%──────────────────────────────────────────────────────────────────────────────
\section{Criterion C7: Non-Trivial Sector (Exclusion of $p = 2$)}
\label{sec:spr:c7}
%──────────────────────────────────────────────────────────────────────────────

The prime $p = 2$ is prime-stable (it satisfies the fixed-point condition
with $B(2) = 1$), but belongs to the \emph{trivial sector} of UBT for
the following independent structural reasons:
\begin{enumerate}
  \item $B(2) = 1$: the one-loop coefficient equals the minimal positive integer.
        The logarithmic term in $V_{\mathrm{eff}}(n) = n^2 - B\,n\ln n$ at $n=2$
        contributes $B \cdot 2 \cdot \ln 2 \approx 1.386$, which is comparable
        in magnitude to the kinetic term $n^2 = 4$.  There is no separation of
        scales between the two contributions, making the effective potential
        structurally degenerate (one regime) rather than exhibiting the
        logarithmic attractor behaviour seen for large $n^* \gg 1$.
  \item $g(X_0(2)) = 0$: the modular curve has genus zero, so the space
        of cusp forms $S_2(\Gamma_0(2))$ is empty and no Hecke eigenvalue
        signal is possible.
  \item The stability interval has no lower bound ($B_{\mathrm{low}}
        = -\infty$), so the lower-margin condition is vacuously satisfied;
        $p = 2$ passes the fixed-point test for a degenerate reason.
  \item Physical: the electromagnetic coupling $\alpha$ corresponds to the
        U(1) sector of the electroweak theory; the UBT winding number
        must be large (of order $\alpha^{-1} \gg 1$) to reproduce the
        weak coupling $\alpha \approx 1/137$.
        $n = 2$ corresponds to a strong coupling ($\alpha_{\mathrm{bare}} = 0.5$),
        inconsistent with any renormalisation-group flow to the observed value.
\end{enumerate}
$p = 2$ is therefore assigned C7 = 0 (excluded from ranking)
and all odd stable primes receive C7 = 1.

%──────────────────────────────────────────────────────────────────────────────
\section{Composite Scoring Table and Ranking}
\label{sec:spr:ranking}
%──────────────────────────────────────────────────────────────────────────────

Table~\ref{tab:spr:scores} collects all six scored criteria for odd stable
primes (C7 excludes $p = 2$).  The total score is $\sum_{k=1}^{6} C_k$.
Each criterion is weighted equally; the score scale is $0$–$3$ per criterion,
giving a maximum of $18$.

\begin{table}[h]
\centering
\caption{Composite UBT criterion scores for the five odd stable primes.
  Columns: C1 = stability margin; C2 = Hecke lepton-ratio; C3 = $p$-adic
  isolation; C4 = theta spectral; C5 = gauge sector; C6 = observed
  $\alpha^{-1}$.  Each score $\in \{0,1,2,3\}$.}
\label{tab:spr:scores}
\begin{tabular}{rcccccccc}
\toprule
$p$ & C1 & C2 & C3 & C4 & C5 & C6 & \textbf{Total} & \textbf{Rank} \\
\midrule
127 & 3 & 0 & 2 & 0 & 0 & 0 & \textbf{5}  & 3rd \\
137 & 2 & 3 & 3 & 3 & 3 & 3 & \textbf{17} & \textbf{1st} \\
139 & 1 & 1 & 1 & 1 & 0 & 2 & \textbf{6}  & 2nd \\
151 & 2 & 0 & 2 & 0 & 0 & 0 & \textbf{4}  & 5th \\
157 & 0 & 0 & 0 & 3 & 2 & 0 & \textbf{5}  & 4th \\
\midrule
2   & \multicolumn{6}{c}{(excluded: trivial sector, C7 = 0)} & — & excl. \\
\bottomrule
\end{tabular}
\end{table}

\begin{resultbox}
\textbf{Ranking result} (no $\alpha$ assumed in positions 1–5):
\[
  p = 137 \;\gg\; p = 139 \;>\; p = 127 \;>\; p = 157 \;>\; p = 151,
  \qquad
  p = 2 \text{ excluded (trivial)}.
\]
\textbf{Score breakdown}:
\[
  17 \;\gg\; 6 \;>\; 5 \;=\; 5 \;>\; 4.
\]
The gap between rank 1 ($17$ points) and rank 2 ($6$ points) is $11$ points —
larger than the total score of any other prime.
\end{resultbox}

\subsection{Tiebreaking: $p = 127$ vs.\ $p = 157$}
\label{sec:spr:tiebreak}

Both $p = 127$ and $p = 157$ score $5$ but have complementary profiles:
\begin{itemize}
  \item $p = 127$ earns its score entirely through C1 (maximum stability,
        $\sigma = 1.194$, the largest of all odd stable primes) and C3
        (Mersenne prime $2^7 - 1$, extreme $2$-adic isolation).
        It is structurally robust but absent from the Hecke, theta, and
        gauge-EW criteria.
  \item $p = 157$ earns its score through C4 (theta spectral, $\equiv 1
        \pmod{4}$, same as $p = 137$) and C5 (partial gauge: $\nu_2 = 2$
        but $\nu_3 = 2$).  However, its stability margin is $\sigma = 0.007$
        (essentially at the boundary of stability) — a critical weakness.
\end{itemize}

Applying the stability margin as the primary physical tiebreaker (a prime
at the edge of stability is susceptible to any perturbation of the
one-loop coefficient $B$): $p = 127$ (rank 3) is placed above $p = 157$
(rank 4).

%──────────────────────────────────────────────────────────────────────────────
\section{Interpretation}
\label{sec:spr:interpretation}
%──────────────────────────────────────────────────────────────────────────────

\subsection{Why $p = 137$ dominates}

The prime $p = 137$ is the unique element of $\mathcal{S}$ that scores
positively in all six independent criteria simultaneously:

\begin{enumerate}
  \item \textbf{C1 (stability)}: $B(137) = 46$ is centrally placed in its
        stability interval with nearly symmetric margins $(\Delta_- \approx
        \Delta_+ \approx 0.56)$, the most symmetric configuration in $\mathcal{S}$.
  \item \textbf{C2 (Hecke)}: Unique prime in $[50, 300]$ for which Hecke
        eigenvalues reproduce both charged-lepton mass ratios at $< 0.1\%$
        error, independently of $\alpha$.
  \item \textbf{C3 ($p$-adic)}: Only odd stable prime with $B(p)$ an
        exact integer ($B(137) = 46$), compatible with integer Kac-Moody
        levels.
  \item \textbf{C4 (theta)}: $137 \equiv 1 \pmod{8}$ places it in both
        $\vartheta_3^2$ and $\vartheta_3^3$ spectra — compatible with all
        spatial dimensionalities of $\BB = \CC \otimes_\RR \HH$.
  \item \textbf{C5 (gauge)}: $\nu_2(\Gamma_0(137)) = 2$, $\nu_3 = 0$ gives
        the pure electroweak sector ($\mathrm{SU}(2) \times \mathrm{U}(1)$)
        required for the electromagnetic coupling $\alpha$.
  \item \textbf{C6 (obs.)}: Nearest stable prime to $\alpha^{-1}_{\mathrm{obs}}
        = 137.036$ by $\delta = 0.036$, an order of magnitude closer than the
        second-nearest ($p = 139$, $\delta = 1.964$).
\end{enumerate}

No other prime in $\mathcal{S}$ satisfies more than three of these criteria.

\subsection{The role of the other stable primes}

\begin{description}
  \item[$p = 139$ (rank 2)] The neighboring prime of $137$ is retained
        as a comparison/control prime in the finite stable-prime set.  Earlier
        Set~B / mirror-sector interpretations have been moved to the
        speculative extension and are not part of the canonical ranking claim.
  \item[$p = 127$ (rank 3)] The most robustly stable prime ($\sigma = 1.194$).
        Its Mersenne structure ($2^7 - 1$) gives extreme 2-adic isolation.
        However, $127 \equiv 7 \pmod{8}$ places it outside both the
        $\vartheta_3^2$ and $\vartheta_3^3$ theta spectra, and its gauge sector
        ($\nu_2 = 0$) is incompatible with the EM coupling.  Its high stability
        score suggests it is a \emph{stable but physically inert} sector of the
        UBT vacuum.
  \item[$p = 157$ (rank 4)] Theta-spectrally equivalent to $137$
        ($\equiv 1 \pmod{4}$ and $\equiv 1 \pmod{8}$) and has partial gauge
        compatibility ($\nu_2 = 2$, $\nu_3 = 2$).  Its critical weakness is
        the vanishing stability margin ($\sigma = 0.007$): the slightest upward
        perturbation of $B$ eliminates its prime-stable status.
  \item[$p = 151$ (rank 5)] Like $p = 127$, it satisfies $\equiv 7
        \pmod{8}$ (absent from all theta spectra) and $\nu_2 = 0$ (no EW
        sector).  Its moderate stability ($\sigma = 0.353$) and 2-adic
        isolation ($v_2 = 3$) prevent it from scoring last but provide no
        positive structural signal for $\alpha$.
  \item[$p = 2$ (excluded)] Trivial sector: genus-zero modular curve,
        degenerate stability condition, and coupling strength $\alpha = 0.5$
        incompatible with the weak electromagnetic coupling.
\end{description}

\subsection{Independence of the ranking from $\alpha$}

The top-ranked prime $p = 137$ achieves its position on the basis of five
criteria (C1–C5) that are entirely independent of $\alpha$:
C1 uses only $V_{\mathrm{eff}}$ stability intervals derived from the modular
index formula $B(p) = (p+1)/3$ \OPEN; C2 uses Hecke eigenvalues of modular
forms at levels $76, 7, 208$ and lepton masses (not $\alpha$); C3 uses
divisibility arithmetic; C4 uses the Fermat--Legendre sum-of-squares theorems;
C5 uses the Legendre symbol $(-1/p)$ and $(-3/p)$.
Only C6 invokes $\alpha$, and it merely confirms what C1–C5 already select.

\begin{resultbox}
\textbf{Summary}: The prime $p = 137$ ranks first among the six finite
prime-stable primes by every independent UBT physical criterion.  Its score
of~$17/18$ is obtained without using $\alpha$ as an input.  The observational
comparison (C6) confirms the structural selection.  The ranking is:
\[
  \boxed{137 \;\succ\; 139 \;\succ\; 127 \;\succ\; 157 \;\succ\; 151 \;\succ\; (2 \text{ excl.})}
\]
\end{resultbox}

\subsection{Open problems generated by this ranking}

\begin{gapbox}
\begin{enumerate}
  \item \textbf{Gap G-Bmod} [\OPEN]: Derive $B(p) = (p+1)/3$ from the UBT
        action $S[\Theta]$.  Until closed, C1, C3, and C5 rest on a working
        hypothesis.
  \item \textbf{Gap G137-Hk} [\OPEN]: Derive why levels $76, 7, 208$ and
        weights $2, 4, 6$ are selected by $\Theta(q,\tau)$
        (C2 \MC).
  \item \textbf{Gap G-139-control} [\OPEN]: Clarify whether the comparison-prime
        role of $p=139$ has any canonical meaning beyond being a neighboring
        stable-prime control.  Mirror-sector interpretations are tracked only
        in the speculative extension.
  \item \textbf{Gap G-157} [\OPEN]: Clarify why $p = 157$ has a nearly
        identical theta and partial gauge profile to $p = 137$ but is
        essentially unstable ($\sigma = 0.007$).
        Is there a physical UV cutoff on the stable set arising from the
        condition $\sigma(p) > \epsilon_{\mathrm{min}}$?
\end{enumerate}
\end{gapbox}

\ifdefined\INCLUDEMODE
\else
  \end{document}
\fi
